\section{Descenso del gradiente}

\subsection{Formulación}

El descenso del gradiente actualiza el punto actual en la dirección
opuesta al gradiente:
\[
x_{t+1} = x_t - \alpha \nabla f(x_t).
\]
El gradiente se calcula con \texttt{torch.autograd.grad}; la actualización
se realiza fuera del grafo de autodiferenciación (\texttt{torch.no\_grad})
para evitar acumular historial innecesario.

\subsection{Aplicación a cada función}
\label{sec:gd_aplicacion}

Partiendo de un punto aleatorio en $[-10,10]^2$ y ejecutando $P=25$
iteraciones, se usó una tasa de aprendizaje distinta por función, ajustada
manualmente para evitar divergencia:

\begin{itemize}
    \item $f_0$: $\alpha=0.1$.
    \item $f_1$ (Ackley): $\alpha=0.01$.
    \item $f_2$ (Himmelblau): $\alpha=0.002$. Esta función crece como un
    polinomio de grado cuatro, por lo que su gradiente crece rápidamente
    lejos de los mínimos: por ejemplo, en $(10,10)$ el gradiente es
    aproximadamente $(4166, 4318)$. Con una tasa mayor, un solo paso puede
    salir del dominio hacia una zona con un gradiente aún mayor y la
    sucesión diverge.
\end{itemize}

% TODO: Insertar aquí la curva de aprendizaje y la trayectoria
% (plot_func_hist) de la corrida usada como ejemplo para cada función.

\subsection{Calibración de hiperparámetros con Optuna}
\label{sec:gd_optuna}

% TODO: Actualizar esta subsección con la calibración actual del notebook.
% Ahora Optuna ajusta solo alpha con 25 iteraciones fijas y minimiza el
% promedio de f al final en los 6 puntos de ajuste, sin penalizaciones. El
% texto y la tabla siguientes corresponden a una versión anterior (alpha e
% iteraciones, con penalización); los valores vigentes están en la tabla
% "Mejores hiperparámetros encontrados con Optuna" del notebook.

Se calibraron $\alpha$ y la cantidad de iteraciones con Optuna
(muestreador TPE, semilla fija), evaluando cada ensayo como el promedio de
$f(x_{\text{final}})$ sobre 6 puntos iniciales fijos compartidos entre
ensayos. Un ensayo se penaliza (score $=10^6$) si el algoritmo diverge o si
la norma del gradiente final supera $10^{-3}$. Con 25 ensayos por función,
los mejores resultados encontrados fueron:

\begin{table}[H]
\centering
\begin{tabular}{lrrr}
\toprule
Función & Mejor $\alpha$ & Mejores iteraciones & Promedio $f(x_{\text{final}})$ \\
\midrule
$f_0$ & $0.456334$   & $80$  & $2.81\times10^{-168}$ \\
$f_1$ & $0.00763728$ & $225$ & $13.30575$ \\
$f_2$ & $0.00552217$ & $510$ & $1.459\times10^{-29}$ \\
\bottomrule
\end{tabular}
\caption{Mejores hiperparámetros de descenso del gradiente encontrados por
Optuna, en 25 ensayos por función.}
\end{table}

Nótese que Optuna pudo explorar cantidades de iteraciones mayores a
$P=25$ (hasta 80, 225 y 510 según la función), por lo que estos valores no
son directamente comparables con las corridas de 25 iteraciones de la
sección anterior. En Ackley el promedio permanece lejos de cero porque los
puntos iniciales fijos caen en cuencas de atracción de distintos mínimos
locales.

% TODO: Documentar el protocolo de 10 corridas (máximo 50 iteraciones), que
% ya está ejecutado en la sección "Evaluación con los mejores
% hiperparámetros" del notebook: tabla de convergencia por corrida, valor
% promedio de la función minimizada e iteraciones promedio para converger.
% TODO: Insertar la trayectoria (curvas de nivel) y la curva de aprendizaje
% de la corrida con convergencia más rápida para cada función (sección
% "Mejores corridas por velocidad de convergencia" del notebook).
% TODO: Comparar formalmente contra RMSProp (punto 2f), ver
% Sección~\ref{sec:comparacion}.
